% GENERATED FILE. Edit canonical lessons, then run npm run generate:books.
% canonical-source-hash: fnv1a64:f31224d58f98fede
% canonical-lessons: FA-C417-shabake, FA-C417-shabake-ye-ejtemai, FA-C417-ettesal, FA-C417-dastresi, FA-C417-majazi

\chapter{Online}
\label{ch:fa-online}
\hlchaptermodality{417}

\begin{chapteropening}
\textbf{By the end of this chapter:} \emph{I can talk about networks and social media, my internet connection, access to the internet and online classes.}

Complete the last lesson of chapter 417: I can talk about networks and social media, my internet connection, access to the internet and online classes.
\end{chapteropening}

\section[\texorpdfstring{\fa{شبکه} (shabake)}{شبکه (shabake)}]{\fa{شبکه} (shabake) — a network; a TV channel}

\label{lesson:FA-C417-shabake}



\begin{quote}
Before the new one: say the Persian for to delete, then the Persian for an update.
\end{quote}

\subsection*{You'll want to know: \fa{شبکه}}
\textbf{\fa{شبکه}} (\emph{shabake}) — \textquotedblleft{}a network; a TV channel\textquotedblright{}.

\begin{quote}
\fa{این} \fa{شبکه} \fa{فیلمهای} \fa{خوبی} \fa{پخش} \fa{میکند}.
\end{quote}

\begin{quote}
\emph{in shabake filmhâ-ye khubi pakhsh mikonad.}
\end{quote}

\begin{quote}
This channel shows good films.
\end{quote}

Plural \emph{shabakehâ}. A computer network too.

\begin{cousinweb}
\textbf{\fa{شبکه}} was borrowed from Arabic \emph{shabaka} \textquotedblleft{}a net\textquotedblright{} (EtymDB): a network is a net of links.
\end{cousinweb}

\begin{grammarlens}[title={What you've built}]
One more word for telling, explaining and describing at B1.
\end{grammarlens}

\subsection*{Your turn}
\begin{itemize}
\raggedright
  \item \emph{Say it:} \emph{shabake}
  \item \emph{Say it:} \emph{shabake}, once more
  \item \emph{Say it:} a sentence of your own with \emph{shabake}
  \item \emph{Recall:} say the Persian for to delete, then the Persian for an update, then say \emph{shabake} again
\end{itemize}

\begin{tcolorbox}[breakable,colback=teal!4,colframe=teal!35!black,title={Before you move on}]
What does \fa{شبکه} mean? (\textquotedblleft{}A network; a TV channel\textquotedblright{}.) Say it once more.
\end{tcolorbox}
\section[\texorpdfstring{\fa{شبکه} \fa{اجتماعی} (shabake-ye ejtemâ'i)}{شبکه اجتماعی (shabake-ye ejtemâ'i)}]{\fa{شبکه} \fa{اجتماعی} (shabake-ye ejtemâ'i) — a social network}

\label{lesson:FA-C417-shabake-ye-ejtemai}



\begin{quote}
Before the new one: say the Persian for an update, then the Persian for a network.
\end{quote}

\subsection*{You'll want to know: \fa{شبکه} \fa{اجتماعی}}
\textbf{\fa{شبکه} \fa{اجتماعی}} (\emph{shabake-ye ejtemâ'i}) — \textquotedblleft{}a social network\textquotedblright{}.

\begin{quote}
\fa{او} \fa{در} \fa{یک} \fa{شبکه} \fa{اجتماعی} \fa{عکس} \fa{میگذارد}.
\end{quote}

\begin{quote}
\emph{u dar yek shabake-ye ejtemâ'i 'aks migozârad.}
\end{quote}

\begin{quote}
She posts photos on a social network.
\end{quote}

\emph{Shabake} \textquotedblleft{}network\textquotedblright{} + ezafe + \emph{ejtemâ'i} \textquotedblleft{}social\textquotedblright{}.

\begin{cousinweb}
\textbf{\fa{شبکه} \fa{اجتماعی}} joins \emph{shabake}, from Arabic \emph{shabaka} \textquotedblleft{}a net\textquotedblright{}, to \emph{ejtemâ'i}, from Arabic \emph{ijtimā'ī} \textquotedblleft{}social\textquotedblright{} (both in EtymDB), with a Persian ezafe.
\end{cousinweb}

\begin{grammarlens}[title={What you've built}]
One more word for telling, explaining and describing at B1.
\end{grammarlens}

\subsection*{Your turn}
\begin{itemize}
\raggedright
  \item \emph{Say it:} \emph{shabake-ye ejtemâ'i}
  \item \emph{Say it:} \emph{shabake-ye ejtemâ'i}, once more
  \item \emph{Say it:} a sentence of your own with \emph{shabake-ye ejtemâ'i}
  \item \emph{Recall:} say the Persian for an update, then the Persian for a network, then say \emph{shabake-ye ejtemâ'i} again
\end{itemize}

\begin{tcolorbox}[breakable,colback=teal!4,colframe=teal!35!black,title={Before you move on}]
What does \fa{شبکه} \fa{اجتماعی} mean? (\textquotedblleft{}A social network\textquotedblright{}.) Say it once more.
\end{tcolorbox}
\section[\texorpdfstring{\fa{اتصال} (ettesâl)}{اتصال (ettesâl)}]{\fa{اتصال} (ettesâl) — a connection}

\label{lesson:FA-C417-ettesal}



\begin{quote}
Before the new one: say the Persian for a network, then the Persian for a social network.
\end{quote}

\subsection*{You'll want to know: \fa{اتصال}}
\textbf{\fa{اتصال}} (\emph{ettesâl}) — \textquotedblleft{}a connection\textquotedblright{}.

\begin{quote}
\fa{اتصال} \fa{اینترنت} \fa{امروز} \fa{خیلی} \fa{کند} \fa{است}.
\end{quote}

\begin{quote}
\emph{ettesâl-e internet emruz kheyli kond ast.}
\end{quote}

\begin{quote}
The internet connection is very slow today.
\end{quote}

\emph{Ettesâl dâshtan} — to be connected.

\begin{cousinweb}
\textbf{\fa{اتصال}} was borrowed from Arabic \emph{ittiṣāl} \textquotedblleft{}connection, contact\textquotedblright{} (EtymDB).
\end{cousinweb}

\begin{grammarlens}[title={What you've built}]
One more word for telling, explaining and describing at B1.
\end{grammarlens}

\subsection*{Your turn}
\begin{itemize}
\raggedright
  \item \emph{Say it:} \emph{ettesâl}
  \item \emph{Say it:} \emph{ettesâl}, once more
  \item \emph{Say it:} a sentence of your own with \emph{ettesâl}
  \item \emph{Recall:} say the Persian for a network, then the Persian for a social network, then say \emph{ettesâl} again
\end{itemize}

\begin{tcolorbox}[breakable,colback=teal!4,colframe=teal!35!black,title={Before you move on}]
What does \fa{اتصال} mean? (\textquotedblleft{}A connection\textquotedblright{}.) Say it once more.
\end{tcolorbox}
\section[\texorpdfstring{\fa{دسترسی} (dastresi)}{دسترسی (dastresi)}]{\fa{دسترسی} (dastresi) — access}

\label{lesson:FA-C417-dastresi}



\begin{quote}
Before the new one: say the Persian for a social network, then the Persian for a connection.
\end{quote}

\subsection*{You'll want to know: \fa{دسترسی}}
\textbf{\fa{دسترسی}} (\emph{dastresi}) — \textquotedblleft{}access\textquotedblright{}.

\begin{quote}
\fa{در} \fa{روستا} \fa{دسترسی} \fa{به} \fa{اینترنت} \fa{کم} \fa{است}.
\end{quote}

\begin{quote}
\emph{dar rustâ dastresi be internet kam ast.}
\end{quote}

\begin{quote}
In the village there is little access to the internet.
\end{quote}

\emph{Dast} \textquotedblleft{}hand\textquotedblright{} + \emph{res-}, from \emph{residan} \textquotedblleft{}to reach\textquotedblright{}, + \emph{-i}. \emph{Dastresi dâshtan be} — to have access to.

\begin{cousinweb}
\textbf{\fa{دسترسی}} is native Persian, built from \emph{dast} \textquotedblleft{}hand\textquotedblright{} + \emph{res-}, the present stem of \emph{residan} \textquotedblleft{}to reach, to arrive\textquotedblright{}, + \emph{-i}. Access is what your hand can reach.
\end{cousinweb}

\begin{grammarlens}[title={What you've built}]
One more word for telling, explaining and describing at B1.
\end{grammarlens}

\subsection*{Your turn}
\begin{itemize}
\raggedright
  \item \emph{Say it:} \emph{dastresi}
  \item \emph{Say it:} \emph{dastresi}, once more
  \item \emph{Say it:} a sentence of your own with \emph{dastresi}
  \item \emph{Recall:} say the Persian for a social network, then the Persian for a connection, then say \emph{dastresi} again
\end{itemize}

\begin{tcolorbox}[breakable,colback=teal!4,colframe=teal!35!black,title={Before you move on}]
What does \fa{دسترسی} mean? (\textquotedblleft{}Access\textquotedblright{}.) Say it once more.
\end{tcolorbox}
\section[\texorpdfstring{\fa{مجازی} (majâzi)}{مجازی (majâzi)}]{\fa{مجازی} (majâzi) — virtual, online}

\label{lesson:FA-C417-majazi}



\begin{quote}
Before the new one: say the Persian for a connection, then the Persian for access.
\end{quote}

\subsection*{You'll want to know: \fa{مجازی}}
\textbf{\fa{مجازی}} (\emph{majâzi}) — \textquotedblleft{}virtual, online\textquotedblright{}.

\begin{quote}
\fa{این} \fa{ترم} \fa{کلاسها} \fa{مجازی} \fa{است}.
\end{quote}

\begin{quote}
\emph{in term kelâshâ majâzi ast.}
\end{quote}

\begin{quote}
This term the classes are online.
\end{quote}

\emph{Fazâ-ye majâzi} — cyberspace, the internet. It also means figurative.

\begin{cousinweb}
\textbf{\fa{مجازی}} is built on \emph{majâz} \textquotedblleft{}figure of speech, a non-literal sense\textquotedblright{}, usually traced to Arabic \emph{majāz}, with the suffix \emph{-i}: not real in the literal sense.
\end{cousinweb}

\begin{grammarlens}[title={What you've built}]
One more word for telling, explaining and describing at B1.
\end{grammarlens}

\subsection*{Your turn}
\begin{itemize}
\raggedright
  \item \emph{Say it:} \emph{majâzi}
  \item \emph{Say it:} \emph{majâzi}, once more
  \item \emph{Say it:} a sentence of your own with \emph{majâzi}
  \item \emph{Recall:} say the Persian for a connection, then the Persian for access, then say \emph{majâzi} again
\end{itemize}

\begin{tcolorbox}[breakable,colback=teal!4,colframe=teal!35!black,title={Before you move on}]
What does \fa{مجازی} mean? (\textquotedblleft{}Virtual, online\textquotedblright{}.) Say it once more.
\end{tcolorbox}
